% 原书第 10 章，印刷页 423--438。
\chapter{大语言模型的演进与前沿}
\label{ch:frontiers}
\chaptercredit{10}
\chapterabstract{本章回顾大语言模型的演进，重点介绍若干趋势和进展。使用合成数据有助于模型继续扩展规模，扩大上下文窗口可提高模型分析信息的能力，改进训练方法则可提高效率和知识迁移的深度。本章也讨论从传统 Transformer 架构转向状态空间模型等替代方案，以改善可扩展性和计算效率。小型模型、技术普及和领域专用模型等趋势，表明人工智能方案正变得更易获取，也更便于针对行业和任务定制。最后介绍大语言模型技术的前沿，以及在逐步取代传统技术的智能体应用和搜索引擎中的用法。}

\section{引言}
大语言模型的演进涉及架构、训练方法和应用趋势。模型架构与训练效率方面的创新，使大语言模型能够完成范围更广、更复杂的任务。使用合成数据和扩大上下文窗口，体现了提高模型能力与表现的持续努力。小型语言模型、借助开源项目推动技术普及、领域专用语言模型等新趋势，也表明应用范围正在扩大，获得这些技术变得更容易。大语言模型智能体和增强搜索等前沿方向，提升了复杂任务执行和信息检索的能力，进一步扩展了大语言模型在各领域的用途。

\section{大语言模型的演进}
\label{sec:ch10-evolution}

\subsection{合成数据}
\label{sec:ch10-synthetic-data}
随着人工智能模型规模增大，易于获取的高质量网络数据日渐耗尽，迫切需要采用合成数据，使模型能继续发展并扩大规模。这一趋势的前提是，增加数据量能提高模型表现，尤其是在复杂、少见的任务上。有观点认为，合成数据反映的是已有数据分布，未必能推动当时最先进的模型继续提高。也有人认为，合成数据的多样性有助于改进模型。

Anthropic 在其人工智能模型中广泛使用合成数据，Claude 2.1 是代表。它可以更准确地拒答无法回答的问题，从而提高稳健性。Anthropic 的宪法式人工智能（Constitutional AI，CAI）方法主要通过两种方式使用合成数据。先依据一组伦理原则批评并改写响应，再生成成对偏好数据进行强化学习。后者属于基于人工智能反馈的强化学习（RLAIF），参见第~5~章。依照原则修正指令响应和依照原则利用反馈进行强化学习的结合，被证明是有效的。虽然团队相对较小，Anthropic 仍在合成数据利用和模型训练方面取得了较好的结果（\refcite{ch10-ref-001}{Bai et al., 2022}）。

Alpaca、Vicuna 等模型使用合成数据对 Llama 模型进行监督微调，提高 70 亿至 130 亿参数规模的模型表现（\refcite{ch10-ref-020}{Peng et al., 2023}; \refcite{ch10-ref-025}{Taori et al., 2023}）。当时常用的方法包括 Self-Instruct，即让大语言模型从种子指令出发生成多样化的指令数据。不过，如何提高数据多样性仍处于探索初期。也有人使用 GPT-4 等模型，把质量较低的网络提示改作训练指令。

UltraFeedback 等合成偏好数据集收集用户编写的提示及模型补全结果，供基于反馈的强化学习训练使用（\refcite{ch10-ref-005}{Cui et al., 2023}）。Teknium1 一直在用合成指令训练基于 Mistral 的 OpenHermes 等模型（\refcite{ch10-ref-008}{Gallego, 2024}）。英特尔当时推出的 Neural-Chat-v3-1 使用 DPO 纳入合成偏好。伯克利的 Starling 模型使用 Nectar，即由 GPT-4 标注的排序数据集。Nectar 汇总来自 GPT-4、GPT-3.5-instruct、GPT-3.5-turbo、Mistral-7B-Instruct 和 Llama-2-7B 等模型的提示及分数，共得到 380 万个成对比较。Starling 在 MT Bench 的 7B 模型中达到了当时的先进水平，不过也有数据污染的疑虑（\refcite{ch10-ref-036}{Zhu et al., 2023a}）。基于人工智能反馈的质量—多样性方法（Quality-Diversity through AI Feedback，QDAIF）用演化算法增加数据多样性（\refcite{ch10-ref-003}{Bradley et al., 2023}）。Evol-Instruct 结合 GPT-4 的反馈与规则系统，生成多样而高质量的指令（\refcite{ch10-ref-032}{Xu et al., 2023}）。

\subsection{更大的上下文窗口}
\label{sec:ch10-long-context}
上下文窗口决定大语言模型在一次处理过程中能利用哪些信息，相当于短期记忆，适用于生成与对话任务。窗口越大，提示可包含越丰富、越详细的示例，从而提高模型从提示中学习的能力，使响应更准确、更相关。较大的窗口还能帮助模型理解相隔很远的文本片段并建立联系，这对详细文档摘要、问答和聊天机器人等任务尤其有用，也有助于长时间交互保持连贯。

GPT 系列模型的上下文窗口逐代扩大。按原书列出的数字，GPT-3 起初上限约为 2,000 个词元，初版 GPT-4 增至 4,096 个词元，GPT-4–32k 变体增至 32,768 个词元。当时的 GPT-4 Turbo 支持最多 128,000 个词元，按原书数字约为初版 GPT-4 的 32 倍、GPT-4–32k 的 4 倍，提升了模型分析大段文本的能力。\footnote{本段保留原书的数值和比较。OpenAI 的 \href{https://developers.openai.com/api/docs/models/gpt-4}{GPT-4 模型规格}列出的上下文窗口是 8,192 个词元。按此规格，128,000 约为 8,192 的 15.6 倍。}Anthropic 的 Claude 支持 9,000 个词元，其后继型号 Claude 2 将容量扩大到 100,000 个词元，可以在一个提示中处理最多约 75,000 个英文词的文档。Meta AI 的 Llama 系列也有支持 100,000 个以上词元的模型。

旋转位置嵌入（RoPE）直接将词元位置嵌入 Transformer 模型（\refcite{ch10-ref-024}{Su et al., 2024}）。该方法根据各词元在序列中的位置施加旋转，在上下文窗口扩大时仍能表征词元之间的相对位置。位置跳跃式训练（Positional Skip-wise Training，PoSE）通过有控制地跳过位置，高效扩展大语言模型的上下文窗口，改善训练和推理阶段处理长上下文的能力（\refcite{ch10-ref-037}{Zhu et al., 2023b}）。LongRoPE 采用改进的旋转嵌入来处理极长的输入，将上下文窗口扩展到超过 200 万个词元（\refcite{ch10-ref-006}{Ding et al., 2024}）。

Munkhdalai et al. (2024) 提出一种名为 Infini-attention 的注意力方法，使大语言模型能够处理极长的输入。它把压缩记忆与局部注意力及长程线性注意力相结合，在上下文检索任务中处理了多达 100 万个词元，在书籍摘要任务中处理了 50 万个词元。

\subsection{训练加速}
本节讨论提高 Transformer 模型效率的若干技术。尽管 Transformer 显著改善了序列建模任务，标准注意力的计算与存储开销仍随序列长度以平方量级增长。参数共享、剪枝、混合精度和微批处理等技术缓解了这些问题，使 Transformer 更易实际采用并推广（\refcite{ch10-ref-007}{Fournier et al., 2023}）。

梯度检查点（gradient checkpointing）在前向传播时选择性保存激活，在反向传播时重新计算其余激活，以节省内存。以额外计算换取内存的做法，使模型可以增加层数，而无需按层数线性增加保存激活的内存。参数共享在网络的不同位置复用同一组参数，从而减少可训练参数的数量。剪枝在训练后去掉较不重要的权重，以提高模型效率。结构化剪枝针对层或注意力头等组件，非结构化剪枝则针对单个权重。剪枝有助于构建体积更小、运行更快、适合现代计算硬件的模型。

为了提高训练速度并降低深度学习模型的内存消耗，现代 GPU 和 TPU 使用混合精度技术。它们以半精度（16 位）执行计算，同时保留单精度的主权重副本，以维持数值稳定性。NVIDIA 的自动混合精度（Automatic Mixed Precision）便于与 TensorFlow、PyTorch 和 MXNet 等框架集成。GPipe 用微批处理把大型模型分布到多个处理单元，便于扩大模型规模并改善性能。它把小批次拆成更小的微批次，使这些微批次可并行经过流水线，并降低前向与反向传播的内存需求。随着加速器数目增加，该方法可扩展模型规模，提高训练吞吐量，同时保持较高的计算效率。

\subsection{多词元生成}
采用常规下一词元预测的大语言模型需要大量资源，也常常难以有效捕捉长程依赖。Meta 的研究提出用多词元预测训练大语言模型。该方法同时预测多个未来词元，以改善效率和表现（\refcite{ch10-ref-009}{Gloeckle et al., 2024}）。在论文的试验中，这项技术可使推理速度达到原来的约 3 倍，并提高样本效率，对较大的模型和编程任务尤其有效。Meta 的 130 亿参数模型在 HumanEval 和 MBPP 基准上的问题求解能力分别提高了 12\% 和 17\%。

该方法用共享的模型主干将输入序列处理为潜在表征，再由多个输出头分别预测不同位置的未来词元。训练时，共享主干使模型能够并行作出多项词元预测，无需为每项预测独立运行主干。推理时，模型可用训练好的输出头在一次迭代中生成多个候选词元，进一步加快生成并减少延迟。

\subsection{知识蒸馏}
\label{sec:ch10-distillation}
知识蒸馏（knowledge distillation，KD）把大型、复杂的教师模型所学的知识迁移给更小、更高效的学生模型。大规模模型的计算需求较高，部署时又面临资源限制，因此知识蒸馏对实际部署很重要。GPT-4、Gemini 等大语言模型兴起后，蒸馏的目标从缩小模型或模仿输出，发展为更细致地迁移深层知识。

大语言模型具有丰富、细致的理解能力，单靠剪枝或量化等传统压缩方法无法完整保留。较新的蒸馏方法用精心设计的提示提取特定知识或能力，利用大语言模型在自然语言处理、推理和问题求解等方面的专长。这样可以更有针对性地迁移知识，并按需要侧重特定技能或领域。

知识蒸馏还试图迁移推理方式、偏好对齐和伦理价值等较抽象的属性。相关技术让学生模型学习教师模型的思考与决策模式。常用手段是思维链提示，让学生模型学习并复现教师模型给出的推理过程，提高处理复杂任务的能力。

Xu et al. (2024) 在综述中将知识蒸馏划分为三个主要方面，即蒸馏算法、技能蒸馏和垂直领域蒸馏，每一方面又涵盖多种方法和子题。

蒸馏算法研究如何从教师模型构建知识，并将其纳入学生模型。它涉及标注、扩充、数据整理、特征理解、反馈机制以及自知识生成。为了有效迁移知识，还使用监督微调、散度最小化、强化学习和排序优化等多种学习方式，使开源模型有可能达到或超过专有模型的能力。

技能蒸馏研究如何通过知识蒸馏提高特定能力，包括利用上下文、遵循指令、检索增强生成、对齐思考方式、角色或偏好建模，以及价值对齐。它也涉及在自然语言理解与生成、信息检索、推荐系统、文本生成评估、代码生成等自然语言处理任务上的专门化，并考察蒸馏如何提高大语言模型处理多模态输入的能力，扩展适用场景。

垂直领域蒸馏考察知识蒸馏在法律、医疗、金融和科学等专业领域的应用，说明它如何让大语言模型适应特定行业需求。这体现了知识蒸馏对领域专用人工智能方案的影响，也说明它可用于满足不同行业的需求。

\subsection{注意力之后的架构}
状态空间模型（state space model，SSM）是深度学习技术演进中的一个研究方向，可用于处理卷积神经网络（CNN）、循环神经网络（RNN）、图神经网络（GNN）乃至 Transformer 等传统架构的局限。状态空间模型用状态变量表征动态系统，其思想源自控制理论和计算神经科学。Mamba 改善了计算效率，论文报告其推理速度约为 Transformer 的 5 倍，计算量可随序列长度线性增长。它使用随输入变化的状态空间模型，更好地根据内容处理信息。在语言、音频和基因组等任务上，其表现显著超过同等规模的 Transformer，并可媲美规模为其两倍的模型（\refcite{ch10-ref-011}{Gu and Dao, 2023}）。

在语言建模中，研究人员探索了用于长程语言建模的门控状态空间（gated state space，GSS）等方法，显著提高速度并减少计算开销（\refcite{ch10-ref-017}{Mehta et al., 2022}）。结构化状态空间序列模型（S4）为状态空间模型引入了更高效的参数化方式，减少计算开销，并在多种基准上表现出色。S4 在序列 CIFAR-10 及图像和语言建模等任务上达到或超过先前模型，生成速度达到所比较方法的 60 倍，并在 Long Range Arena 基准上创下当时的新纪录，可处理最长 16,000 的序列（\refcite{ch10-ref-012}{Gu et al., 2021}）。

\section{大语言模型的发展趋势}
\label{sec:ch10-trends}

\subsection{小型语言模型}
大语言模型推动了许多领域的发展，但所需的大量计算资源通常只有资源充足的组织才能承担。越来越多研究者尝试用小得多的模型复现大型模型的能力。小型语言模型（small language model，SLM）是大语言模型的缩小版，参数量在数百万至数十亿之间，远少于某些大语言模型的数千亿或数万亿参数。较小的规模带来以下好处。

\begin{enumerate}
\item \textbf{效率。}小型语言模型耗电量更小、内存需求更低，适合部署在较小的设备上，例如边缘计算场景。这便于开发端侧聊天机器人和个人移动助手等应用，直接在用户设备上运行。
\item \textbf{可获取性。}资源需求较低，使更多开发人员和组织能够采用小型语言模型。即使小团队和独立研究人员，也可在没有大量基础设施的情况下利用语言模型，从而推动人工智能技术普及。
\item \textbf{定制。}小型语言模型更易适配特定领域和任务，可开发针对具体需要的专用模型。在细分应用中，这种定制可以提高表现和准确率。
\item \textbf{安全与隐私。}较小的模型可部署在本地或私有云环境中，减少数据泄露的可能。这对金融和医疗等处理敏感信息的行业尤其重要，因为这些行业需要保持对数据的控制。
\end{enumerate}

以下是原书选列的几种当时受到业界关注的小型语言模型。
\begin{enumerate}
\item \textbf{Llama。}原书将 Llama-3 描述为可开放获取的 27 亿参数模型，擅长处理细致的语言任务、翻译和对话生成（\refcite{ch10-ref-028}{Touvron et al., 2023}）。\footnote{此处保留原书表述供对照。\href{https://ai.meta.com/blog/meta-llama-3/}{Meta 于 2024 年 4 月公布的 Llama 3}有 80 亿和 700 亿参数两种规模。原书所引 Touvron et al. (2023) 的文献介绍的是 Llama 2。}
\item \textbf{Phi-2。}微软开发的 27 亿参数模型，在数学推理、常识评估和逻辑任务上表现突出（\refcite{ch10-ref-013}{Javaheripi et al., 2023}）。Phi-2 用合成数据训练，在阅读理解和文本摘要等任务上可与十倍规模的模型竞争，有时还能超过它们。
\item \textbf{Mistral 7B。}该模型有 73 亿个参数，表现优于此前的 Llama 模型，且接近专用代码模型（\refcite{ch10-ref-014}{Jiang et al., 2023a}）。它使用分组查询注意力加快处理，并用滑动窗口注意力高效处理较长的文本序列。
\item \textbf{Gemma 2B 和 Gemma 7B。}这两种规格均有预训练版和指令调优版，在文本任务上表现突出，在 18 项评估中有 11 项优于规模相近的开放模型（\refcite{ch10-ref-027}{Team et al., 2024}）。Gemma 的开发也重视安全和责任，以提高实际应用中的可靠性。
\item \textbf{Vicuna-13B。}这是一个基于 Llama-13B 的开源对话模型，用用户共享的对话微调。初步评估用 GPT-4 作裁判，结果表明 Vicuna-13B 的质量得分超过 ChatGPT 和 Google Bard 的 90\%，并在大部分测试中优于 Llama、Alpaca 等模型（\refcite{ch10-ref-020}{Peng et al., 2023}）。
\end{enumerate}

\subsection{技术普及}
在原书写作前的数月里，开源社区影响力不断扩大，推动了大语言模型的多项变化。开源强调共同开发、透明性与免费获取，对大语言模型的发展影响深远。相关开源项目涉及预训练数据、模型与架构、指令调优数据集、对齐调优数据集，甚至硬件。

Petals 为没有高端硬件、难以使用 BLOOM-176B 和 OPT-175B 等大模型的研究人员提供了一种途径（\refcite{ch10-ref-002}{Borzunov et al., 2022}）。它汇集参与者愿意共享的 GPU 算力，协同执行大模型推理和微调。在交互应用中，它比卸载到主内存的方式更快，可以在消费级 GPU 上以大约每秒一步的速度推理。

Hugging Face 的 ZeroGPU 项目在原书所述的时期使用 NVIDIA A100 GPU，通过 Spaces 应用提供按需共享的 GPU 资源，旨在让更多人能够获得算力，并减少小型组织的开销。

预训练、指令调优和对齐调优等用途的数据集持续开放给社区。贡献者经常在网上发布开源数据集，LLMDataHub 和 Open LLM Datasets 等项目集中整理这些资源，方便从事大语言模型开发的人员和研究者查找和使用。

OpenLLM 让开发人员通过与 OpenAI 兼容的 API 端点在本地或云端运行 Llama-2、Mistral 等开源大语言模型（\refcite{ch10-ref-021}{Pham et al., 2023}）。该平台支持多种大语言模型，方便应用切换 API，并用 BentoML 优化模型服务和简化云端部署。

第~8~章详细讨论了开源大语言模型。读者还可查阅当时的 Hugging Face 排行榜，了解模型的发布和排名。

\subsection{领域专用语言模型}
领域专用语言模型（domain-specific language model，DSLM）针对特定行业或领域设计，弥补通用模型的局限。这些模型用领域专用的数据和术语进行微调，适合需要准确性的复杂或受监管环境。定向训练让模型能给出更准确、符合语境的响应，减少通用模型遇到专业内容时可能出现的错误与幻觉。

律师、医疗服务提供者和金融分析师等专业人员依赖准确、可靠的信息，因而可以从领域专用语言模型中获益。模型聚焦较窄的范围，吸收行业术语，以便处理指定领域的工作流程。随着企业逐渐认识到定制人工智能方案的价值，原书预计，到 2027 年，企业采用的生成式人工智能模型中将有一半以上是面向特定行业或职能的领域专用模型。

法律领域的 SaulLM-7B 由 Equall.ai 开发，是用法律领域的预训练与微调来处理复杂法律语言、改善法律应用任务表现的例子（\refcite{ch10-ref-004}{Colombo et al., 2024}）。医疗领域开发了 GatorTron、Codex-Med、Galactica 和 Flan-PaLM 等模型，以处理医疗数据和临床信息中的细节，探索人工智能在诊断与患者照护管理方面的能力（\refcite{ch10-ref-023}{Singhal et al., 2023}; \refcite{ch10-ref-026}{Taylor et al., 2022}; \refcite{ch10-ref-034}{Yang et al., 2022}, 2023）。金融领域的 BloombergGPT 和 FinBERT 等模型则用大量金融数据训练，以改善风险管理和金融分析等任务（\refcite{ch10-ref-016}{Liu et al., 2021}; \refcite{ch10-ref-031}{Wu et al., 2023}）。

% 原书第 10.4 和 10.5 节。
\section{新的前沿}
\label{sec:ch10-new-frontiers}

\subsection{大语言模型智能体}
\label{sec:ch10-agents}
大语言模型智能体是一种利用大语言模型完成复杂任务的框架。此类智能体是模块化程序，可以读取用户请求，推理完成请求所需的步骤，创建子任务并分配给各模块，最后综合结果并返回给用户。其特点是把传统的计算逻辑和系统工具与大语言模型结合，通过提示引导模型使用这些工具。

例如，设想要对市售耳机进行市场调研，再根据价格、功能和用户评价选出几款。只凭预训练时接触的数据，原书讨论的基础 ChatGPT 无法完成这项任务。大语言模型智能体则可以读取请求，判断需要从网上查找信息，构造相关查询，执行搜索，下载结果网页，处理信息，最后给出若干建议。

不同类型的智能体一般都有若干共同模块，见图~\ref{fig:ch10-agent}。

\begin{figure}[!htbp]
\centering
\begin{tikzpicture}[
  font=\sffamily\small,
  module/.style={draw=blue!50!black,rounded corners=3pt,
    text width=2.35cm,minimum height=.67cm,align=center},
  tool/.style={module,text width=1.55cm,minimum height=.65cm,fill=green!14},
  exchange/.style={{Stealth[length=1.6mm]}-{Stealth[length=1.6mm]}},
  request/.style={-{Stealth[length=1.6mm]}}]
\node[module,fill=yellow!18] (core) at (0,0) {智能体核心};
\node[module,fill=blue!12] (plan) at (-3.85,0) {规划};
\node[module,fill=orange!16] (memory) at (3.85,0) {记忆};
\node[module,fill=black!11] (user) at (0,-1.72) {用户};
\node[tool] (web) at (-3.30,1.72) {网页搜索};
\node[tool] (calc) at (-1.10,1.72) {计算器};
\node[tool] (kb) at (1.10,1.72) {知识库};
\node[tool] (other) at (3.30,1.72) {其他工具};
\draw[exchange] (plan) -- (core);
\draw[exchange] (memory) -- (core);
\foreach \toolname in {web,calc,kb,other}{
  \draw[exchange] (core) -- (\toolname.south);
}
\draw[request] ([xshift=-.48cm]user.north) --
  node[left,font=\sffamily\scriptsize]{任务说明} ([xshift=-.48cm]core.south);
\draw[request] ([xshift=.48cm]core.south) --
  node[right,font=\sffamily\scriptsize]{汇总结果} ([xshift=.48cm]user.north);
\node[below,align=left,font=\sffamily\scriptsize] at (plan.south)
  {分解任务\\推理（思维链、Reflexion）};
\node[below,align=left,font=\sffamily\scriptsize] at (memory.south) {短期记忆\\长期记忆};
\draw (-4.17,2.25) -- (-.45,2.25) (0.45,2.25) -- (4.17,2.25);
\draw (-4.17,2.13) -- (-4.17,2.37) (4.17,2.13) -- (4.17,2.37);
\node at (0,2.25) {工具};
\end{tikzpicture}
\caption{大语言模型智能体的主要模块。核心接收用户任务，从记忆模块取得相关信息，交由规划模块分解子任务，再使用可用工具完成子任务。核心最后综合信息并将结果返回用户。}
\label{fig:ch10-agent}
\end{figure}


\begin{keypoint}
\begin{itemize}
\item \textbf{核心。}定义智能体的特性，接收用户指令，并向其他模块分配任务。这通过一个较长、详细的提示模板实现，用于指导大语言模型如何执行任务。
\item \textbf{规划。}确定完成整体任务所需的一系列步骤。采用思维链等推理方法（见第~4.5.4~节）以及 Reflexion 等迭代改进算法（\refcite{ch10-ref-022}{Shinn et al., 2023}），制定并完善一份由多个子任务组成的计划，再由智能体的各项功能执行。
\item \textbf{工具。}为智能体提供超出大语言模型本身能力的一系列功能。工具可以有很多种，例如网页搜索命令、代码编译器、计算器和各种 API 调用。
\item \textbf{知识。}智能体可在需要时查询的知识库。它可以类似第~7~章介绍的 RAG 系统，也可以是结构化数据库，由语言模型生成 SQL 等调用来查询。
\item \textbf{记忆。}记录从用户与智能体的交互中得到的信息，必要时供任务查阅。有时可分为短期记忆与长期记忆。短期记忆保留当前会话所有交互的细节，长期记忆则保存从多次交互中选取的相关信息。
\end{itemize}
\end{keypoint}

智能体可以针对科学写作（\refcite{ch10-ref-038}{Ziems et al., 2024}）、玩电子游戏（\refcite{ch10-ref-030}{Wang et al., 2023}）、操纵机器人（\refcite{ch10-ref-018}{Michael Ahn, 2022}）等具体任务设计。研究人员也开发了通用智能体，尝试执行用户提出的任意任务。早期例子是 AutoGPT。\footnote{\url{https://github.com/Significant-Gravitas/AutoGPT}}其结构与图~\ref{fig:ch10-agent}相近。它接收用户指令，用核心模块中设计好的提示模板组织工作流，通过思维链和自我批评拟定计划，再利用记忆模块与工具完成任务。它不接受用户对计划的反馈，而是自主尝试整个任务求解过程。另一种当时常用的智能体框架是 BabyAGI。\footnote{\url{https://github.com/yoheinakajima/babyagi}}其整体布局与 AutoGPT 相似，但每完成一个任务就会调整计划，而非依次执行预先确定的一串任务。

\begin{practicaltip}
许多智能体建立在为复杂大语言模型应用设计的开源软件包上。LangChain 和 LlamaIndex 是常见例子，见第~\ref{sec:ch08-frameworks}~节。这些软件包提供调用大语言模型的功能、常用工具的集成、适用于多种场景的提示模板以及 Web 托管功能。BabyAGI 的工作流使用了 LangChain 的集成组件。这种集成也是双向的，LangChain 将 AutoGPT 和 BabyAGI 纳入其产品，智能体系统可以方便地使用 LangChain 中已有的各种模型、向量索引和工具。
\end{practicaltip}

\subsection{大语言模型增强搜索}
\label{sec:ch10-search}
大语言模型的另一应用前沿，是增强不同环境下的搜索能力。大语言模型可用多种方式改进网页与文档搜索算法，大致分为以下几类。
\begin{itemize}
\item 改进排名靠前的搜索结果。
\item 改写或扩展查询。
\item 根据搜索结果进行推理。
\end{itemize}

\subsubsection*{改进排名靠前的搜索结果}
传统搜索根据查询做关键词匹配。它的一个主要限制是，无法返回主题相近却未包含相同关键词的文本。词向量扩大了单个词的搜索范围，使语义相近的词也能匹配。语言模型的句子嵌入进一步匹配语义相近的较长短语。随着嵌入模型的改善，检索有望更准确地限定相关范围，可用第~\ref{sec:ch07-reranking}~节介绍的重排序方法。先用高效搜索算法收集一组排名靠前的结果，再用较慢但能力更强的排序算法动态重排。以网页搜索为例，传统的关键词搜索可以返回几十万条匹配结果，初步排序较为粗略。可根据查询与网页内容的语义相似度，对前 $k$ 条结果重新排序。从 Transformer 出现以来，这类技术已在网页搜索行业普遍使用。能细致分析人类语言含义的大语言模型智能体，又扩展了这类方法的能力。

\subsubsection*{改进查询}
大语言模型还可用于改进查询。第~7~章在 RAG 背景下讨论过其中一些方法。
\begin{itemize}
\item \textbf{查询改写。}训练模型接收人的搜索查询，将其改写为更可能检索到相关结果的形式。
\item \textbf{查询到文档。}生成与所需搜索结果形式相似的合成文档，使之与目标文档更易匹配。
\item \textbf{查询到 SQL。}用大语言模型将针对结构化数据库的自然语言查询转换为代码调用，参见第~\ref{sec:ch07-structured}~节的 SQL 示例。\footnote{原书此处把该示例称为清单 7.4.1.1。第~7.4.1.1~节内包含相应示例。}
\end{itemize}

前瞻式主动检索增强生成（Forward-Looking Active REtrieval augmented generation，FLARE）也是提取更相关信息的一种技术（\refcite{ch10-ref-015}{Jiang et al., 2023b}）。在 FLARE 中，大语言模型从初始查询出发，设想可能包含相关信息的新上下文，接连生成后续查询，执行搜索并纳入新结果。

\subsubsection*{根据搜索结果进行推理}
大语言模型改变搜索的另一种方式，可类比第~7~章所述的检索增强生成。RAG 流水线本就包含对文档的语义搜索。基于 RAG 的搜索引擎可以只返回语义最相近的文档，供用户查阅。这种方法适用于在单份文档或网站域名内搜索，可以改善传统关键词搜索难以识别与查询相近但用词不完全相同的主题这一问题。

先搜索、再生成描述的 RAG 式方法也可用于网页搜索。大语言模型增强的搜索引擎可以在预先建立索引的网页中查找语义最接近的结果，提取页面信息并用作上下文，直接回答查询。You.com 创建了具有类似功能的网页入口。按原书当时的描述，You.com 采用与其他聊天机器人相似的网页查询界面，并与 Google 搜索结合。收到查询后，You.com 在网上搜索与问题有关的信息，处理网页的细节，再以这些上下文回答查询。由于能访问网上新近发布的信息，这类聊天机器人可将互联网作为资料库来做 RAG 式搜索。它可用于综合近期新闻中的信息、市场调研和一般问答。

这类应用仍有其他聊天机器人面临的问题，尤其是模型幻觉和信息来源的可靠性。一个值得注意的例子是，谷歌在原书写作时期开展人工智能新闻聚合与摘要试验，曾对简单问题给出十分不准确的回答，有时甚至让用户吃石头，或把胶水加到比萨上（\refcite{ch10-ref-010}{Grant, 2024}）。出现这些响应，是因为带有讽刺意味的评论或网页被用作上下文。由此可以看到，如果把互联网搜索结果等未经筛选的数据当成大语言模型问答的事实知识库，就可能发生错误。搜索并作答的系统一旦输出不准确的信息，可能造成伤害和声誉损失，因此开发时需要留意信息来源。

\section{结语}
\label{sec:ch10-closing}
在最后一章，我们试着把自己所看到的大语言模型创新前沿介绍给读者。更确切地说，这是多模态大语言模型的创新前沿。正如第~9~章所述，多模态正迅速成为语言建模的新范式。我们对前沿的看法来自撰写本书时查阅的大量文献，未来的面貌却难以预料。人类的创造与发明能力，始终是其中难以预测的变量。毕竟，当 Vaswani 等人发表奠基论文 \refcite{ch10-ref-029}{\textit{Attention Is All You Need}} 时，又有谁能预见 ChatGPT 后来展现的能力？

人类创新的不可预见性，使我们不敢对这一领域的未来作过于具体的判断。人类创新仍是一个影响因素，这本身也反映出大语言模型及其能力在何种程度上接近人类智能。我们要创造能与人类个体或群体同样胜任任务，甚至做得更好的机器，仍有许多东西要学习与理解，还会经历失败，也会取得更多进展。

我们希望，本书的内容和汇集的研究，能为有技能、兴趣、才智或机会参与这段人类技术发展进程的人提供参考。更重要的是，我们希望他们在推进技术时承担责任、遵循伦理，并充分尊重人的尊严。技术都有两面性，可能在众多任务中超越人类能力的技术尤其如此。

愿读者喜欢这段旅程，也希望本书能帮助你理解和掌握大语言模型的研究与应用。人工智能仍在发展，本书也需要随之更新。我们期待将来的版本，也希望读者同样期待。

% PDF 461--463。
\begin{chapterreferences}
\refentry{ch10-ref-001}{Yuntao Bai, Saurav Kadavath, Sandipan Kundu, Amanda Askell, Jackson Kernion, Andy Jones, Anna Chen, Anna Goldie, Azalia Mirhoseini, Cameron McKinnon, et al. Constitutional ai: Harmlessness from ai feedback. \textit{arXiv preprint} \textit{arXiv:2212.08073}, 2022.}
\refentry{ch10-ref-002}{Alexander Borzunov et al. Petals: Collaborative inference and fine-tuning of large models. \textit{arXiv preprint arXiv:2209.01188}, 2022. URL \url{https://arxiv.org/abs/2209.01188}.}
\refentry{ch10-ref-003}{Herbie Bradley, Andrew Dai, Hannah Teufel, Jenny Zhang, Koen Oostermeijer, Marco Bellagente, Jeff Clune, Kenneth Stanley, Grégory Schott, and Joel Lehman. Quality-diversity through ai feedback. \textit{arXiv preprint arXiv:2310.13032}, 2023.}
\refentry{ch10-ref-004}{Pierre Colombo et al. Saullm-7b: A pioneering large language model for law. \textit{arXiv} \textit{preprint arXiv:2403.03883}, 2024.}
\refentry{ch10-ref-005}{Ganqu Cui, Lifan Yuan, Ning Ding, Guanming Yao, Wei Zhu, Yuan Ni, Guotong Xie, Zhiyuan Liu, and Maosong Sun. Ultrafeedback: Boosting language models with high-quality feedback. \textit{arXiv preprint arXiv:2310.01377}, 2023.}
\refentry{ch10-ref-006}{Yiran Ding, Li Lyna Zhang, Chengruidong Zhang, Yuanyuan Xu, Ning Shang, Jiahang Xu, Fan Yang, and Mao Yang. Longrope: Extending llm context window beyond 2 million tokens. \textit{arXiv preprint arXiv:2402.13753}, 2024.}
\refentry{ch10-ref-007}{Quentin Fournier, Gaétan Marceau Caron, and Daniel Aloise. A practical survey on faster and lighter transformers. \textit{ACM Computing Surveys}, 55(14s):1–40, 2023.}
\refentry{ch10-ref-008}{Victor Gallego. Configurable safety tuning of language models with synthetic preference data. \textit{arXiv preprint arXiv:2404.00495}, 2024.}
\refentry{ch10-ref-009}{Fabian Gloeckle, Badr Youbi Idrissi, Baptiste Rozière, David Lopez-Paz, and Gabriel Synnaeve. Better \& faster large language models via multi-token prediction. \textit{arXiv} \textit{preprint arXiv:2404.19737}, 2024.}
\refentry{ch10-ref-010}{Nico Grant. Google’s a.i. search errors cause a furor online. \textit{The New York} \textit{Times}, 2024. URL \url{https://www.nytimes.com/2024/05/24/technology/google-ai-overview-search.html}.}
\refentry{ch10-ref-011}{Albert Gu and Tri Dao. Mamba: Linear-time sequence modeling with selective state spaces. \textit{arXiv preprint arXiv:2312.00752}, 2023.}
\refentry{ch10-ref-012}{Albert Gu, Karan Goel, and Christopher Ré. Efficiently modeling long sequences with structured state spaces. \textit{arXiv preprint arXiv:2111.00396}, 2021.}
\refentry{ch10-ref-013}{Mojan Javaheripi et al. Phi-2: The surprising power of small language models. \textit{Microsoft Research Blog}, 2023.}
\refentry{ch10-ref-014}{Albert Q Jiang et al. Mistral 7b. \textit{arXiv preprint arXiv:2310.06825}, 2023a.}
\refentry{ch10-ref-015}{Zhengbao Jiang, Frank F. Xu, Luyu Gao, Zhiqing Sun, Qian Liu, Jane Dwivedi-Yu, Yiming Yang, Jamie Callan, and Graham Neubig. Active retrieval augmented generation, 2023b.}
\refentry{ch10-ref-016}{Zhuang Liu et al. Finbert: A pre-trained financial language representation model for financial text mining. In \textit{Proceedings of the twenty-ninth international conference on international joint conferences on artificial intelligence}, pages 4513– 4519, 2021.}
\refentry{ch10-ref-017}{Harsh Mehta, Ankit Gupta, Ashok Cutkosky, and Behnam Neyshabur. Long range language modeling via gated state spaces. \textit{arXiv preprint arXiv:2206.13947}, 2022.}
\refentry{ch10-ref-018}{Noah Brown et al. Michael Ahn, Anthony Brohan. Do as i can, not as i say: Grounding language in robotic affordances, 2022.}
\refentry{ch10-ref-019}{Tsendsuren Munkhdalai, Manaal Faruqui, and Siddharth Gopal. Leave no context behind: Efficient infinite context transformers with infini-attention. \textit{arXiv preprint} \textit{arXiv:2404.07143}, 2024.}
\refentry{ch10-ref-020}{Baolin Peng, Chunyuan Li, Pengcheng He, Michel Galley, and Jianfeng Gao. Instruction tuning with gpt-4. \textit{arXiv preprint arXiv:2304.03277}, 2023.}
\refentry{ch10-ref-021}{Aaron Pham et al. OpenLLM: Operating LLMs in production, June 2023. URL \url{https://github.com/bentoml/OpenLLM}.}
\refentry{ch10-ref-022}{Noah Shinn, Federico Cassano, Edward Berman, Ashwin Gopinath, Karthik Narasimhan, and Shunyu Yao. Reflexion: Language agents with verbal reinforcement learning, 2023.}
\refentry{ch10-ref-023}{Karan Singhal et al. Large language models encode clinical knowledge. \textit{Nature}, 620 (7972):172–180, 2023.}
\refentry{ch10-ref-024}{Jianlin Su, Murtadha Ahmed, Yu Lu, Shengfeng Pan, Wen Bo, and Yunfeng Liu. Roformer: Enhanced transformer with rotary position embedding. \textit{Neurocomputing}, 568:127063, 2024.}
\refentry{ch10-ref-025}{Rohan Taori, Ishaan Gulrajani, Tianyi Zhang, Yann Dubois, Xuechen Li, Carlos Guestrin, Percy Liang, and Tatsunori B Hashimoto. Alpaca: A strong, replicable instruction-following model. \textit{Stanford Center for Research on Foundation Models}. \url{https://crfm.stanford.edu/2023/03/13/alpaca.html}, 3(6):7, 2023.}
\refentry{ch10-ref-026}{Ross Taylor et al. Galactica: A large language model for science. \textit{arXiv preprint} \textit{arXiv:2211.09085}, 2022.}
\refentry{ch10-ref-027}{Gemma Team, Mesnard, et al. Gemma: Open models based on gemini research and technology. \textit{arXiv preprint arXiv:2403.08295}, 2024.}
\refentry{ch10-ref-028}{Hugo Touvron et al. Llama 2: Open foundation and fine-tuned chat models, 2023.}
\refentry{ch10-ref-029}{Ashish Vaswani, Noam Shazeer, Niki Parmar, Jakob Uszkoreit, Llion Jones, Aidan N Gomez, Łukasz Kaiser, and Illia Polosukhin. Attention is all you need. In I. Guyon, U. Von Luxburg, S. Bengio, H. Wallach, R. Fergus, S. Vishwanathan, and R. Garnett, editors, \textit{Advances in Neural Information Processing Systems}, volume 30. Curran Associates, Inc., 2017. URL \url{https://proceedings.neurips.cc/paper_files/paper/2017/file/3f5ee243547dee91fbd053c1c4a845aa-Paper.pdf}.}
\refentry{ch10-ref-030}{Zihao Wang, Shaofei Cai, Guanzhou Chen, Anji Liu, Xiaojian Ma, and Yitao Liang. Describe, explain, plan and select: Interactive planning with large language models enables open-world multi-task agents, 2023.}
\refentry{ch10-ref-031}{Shijie Wu, Ozan Irsoy, Steven Lu, Vadim Dabravolski, Mark Dredze, Sebastian Gehrmann, Prabhanjan Kambadur, David Rosenberg, and Gideon Mann. Bloomberggpt: A large language model for finance, 2023.}
\refentry{ch10-ref-032}{Can Xu, Qingfeng Sun, Kai Zheng, Xiubo Geng, Pu Zhao, Jiazhan Feng, Chongyang Tao, and Daxin Jiang. Wizardlm: Empowering large language models to follow complex instructions. \textit{arXiv preprint arXiv:2304.12244}, 2023.}
\refentry{ch10-ref-033}{Xiaohan Xu, Ming Li, Chongyang Tao, Tao Shen, Reynold Cheng, Jinyang Li, Can Xu, Dacheng Tao, and Tianyi Zhou. A survey on knowledge distillation of large language models. \textit{arXiv preprint arXiv:2402.13116}, 2024.}
\refentry{ch10-ref-034}{Xi Yang et al. A large language model for electronic health records. \textit{NPJ digital} \textit{medicine}, 5(1):194, 2022.}
\refentry{ch10-ref-035}{Zhichao Yang et al. Surpassing gpt-4 medical coding with a two-stage approach. \textit{arXiv preprint arXiv:2311.13735}, 2023.}
\refentry{ch10-ref-036}{Banghua Zhu, Evan Frick, Tianhao Wu, Hanlin Zhu, and Jiantao Jiao. Starling-7b: Improving llm helpfulness \& harmlessness with rlaif, 2023a.}
\refentry{ch10-ref-037}{Dawei Zhu, Nan Yang, Liang Wang, Yifan Song, Wenhao Wu, Furu Wei, and Sujian Li. Pose: Efficient context window extension of llms via positional skip-wise training. \textit{arXiv preprint arXiv:2309.10400}, 2023b.}
\refentry{ch10-ref-038}{Caleb Ziems, William Held, Omar Shaikh, Jiaao Chen, Zhehao Zhang, and Diyi Yang. Can large language models transform computational social science?, 2024.}
\end{chapterreferences}

